\documentclass[11pt,letterpaper]{article}
\usepackage[margin=0.8in]{geometry}
\usepackage[T1]{fontenc}
\usepackage{lmodern,microtype,amsmath,amssymb,booktabs,tabularx,graphicx,caption}
\usepackage[hidelinks]{hyperref}
\setlength{\parindent}{0pt}\setlength{\parskip}{5pt}
\setlength{\emergencystretch}{2em}
\setlength{\textfloatsep}{8pt}\setlength{\floatsep}{7pt}
\captionsetup{font=small,labelfont=bf,skip=4pt}
\graphicspath{{results/report_figures/}}
\newcommand{\FullShuttleCost}{597.356}
\newcommand{\OptimalEmptyCost}{576.670}
\newcommand{\WaitCostTwelve}{672.892}
\newcommand{\DispatchCostTwelve}{676.670}
\newcommand{\WaitCostThirteen}{680.770}
\newcommand{\DispatchCostThirteen}{676.670}

\newcommand{\shuttlefigure}[2]{%
\begin{figure}[!ht]\centering
\includegraphics[width=.9\linewidth,height=.20\textheight,keepaspectratio]{#1}
\caption{#2}
\end{figure}}
\begin{document}
\begin{center}{\Large\bfseries Assignment 1: Shuttle Dispatch}\\[3pt]
Value and Policy Iteration\end{center}
\textbf{GitHub username:} FailedMath3U\quad\textbf{Repository:} MLDS\\
\url{https://github.com/FailedMath3U/MLDS}

\section{MDP and modeling assumptions}
At decision period $t$, the state $s_t\in\{0,\ldots,200\}$ is the number waiting before dispatch. The action $a_t\in\{0,1\}$ means wait or dispatch one shuttle. Its capacity is 15. Let $N_t$ be the subsequent arrivals, independently uniform on $\{1,2,3,4,5\}$. Define the passengers remaining after service as $q(s,a)=\max\{s-15a,0\}$. The cost and next state are
\[
c(s,a)=100a+2q(s,a),\qquad
s_{t+1}=\min\{q(s_t,a_t)+N_t,200\}.
\]
Thus dispatch precedes holding cost, and new arrivals enter afterward. Queue length is sufficient for the Markov property under independent arrivals and unrestricted shuttle availability.

The transition probability $P(s'\mid s,a)$ is the probability of next queue $s'$ given state $s$ and action $a$:
\[
P(s'\mid s,a)=\frac15\sum_{k=1}^{5}
\mathbb{I}\!\left\{s'=\min\{q(s,a)+k,200\}\right\}.
\]
$\mathbb{I}\{\cdot\}$ equals one when its condition holds and zero otherwise. Arrival outcomes that reach the same capped state have their probabilities added.

\textbf{Additional assumptions.} At most one shuttle departs each period; there is no fleet-size or travel-time constraint. Overflow above 200 is discarded without penalty. New arrivals first incur holding cost at the next decision. The finite-horizon problem has 500 decisions ($t=0,\ldots,499$) and no terminal cost. Ties choose waiting.

\textbf{Objective and notation.} With discount $\gamma=0.95$, minimize expected discounted cost. $J_t(s)$ is the minimum remaining cost from state $s$ at time $t$, discounting relative to that time:
\[
J_t(s)=\min_\pi\mathbb E_\pi\!\left[
\sum_{u=t}^{499}\gamma^{u-t}c(s_u,a_u)\,\middle|\,s_t=s\right].
\]
Here $u$ counts future periods, $\pi$ is a rule for choosing actions, and $\mathbb E_\pi$ averages arrival randomness under that rule. For an infinite horizon, $J^*(s)$ denotes the optimal cost with no deadline. Plots show cost values: lower is better. If rewards are defined as $-c$, the corresponding reward values are $-J$.

\textbf{Run the code.} From \texttt{assignment1/mlds\_assignment1\_shuttle/}:
\begin{verbatim}
python -m pip install -r requirements.txt
python shuttle.py
jupyter lab shuttle_plots.ipynb
\end{verbatim}
The solver runs all three methods and 14 validation experiments. Evidence is in \texttt{results/summary.json}; \texttt{state\_results.csv} lists all states, and \texttt{shuttle\_results.npz} saves every finite-horizon value/policy. The notebook regenerates plots. Section~\ref{sec:validation} explains each validation experiment.

\clearpage
\section{Part (a): enumeration, $T=500$}
We enumerate every state, action and arrival outcome by backward induction. Set $J_{500}(s)=0$; then for $t=499,\ldots,0$ compute
\[
Q_t(s,a)=c(s,a)+\frac{\gamma}{5}\sum_{k=1}^{5}
J_{t+1}\!\left(\min\{q(s,a)+k,200\}\right),\qquad
J_t(s)=\min_{a\in\{0,1\}}Q_t(s,a).
\]
$Q_t(s,a)$ is the cost of action $a$ now followed by optimal decisions. The minimizing action defines the time-dependent policy $\pi_t(s)$. All 501 value vectors and 500 policy vectors are saved. This reuses continuation values rather than listing exponentially many complete histories; an independently enumerated small tree validates the recursion (Section~\ref{sec:validation}).
\shuttlefigure{a_finite_horizon_value.pdf}{Part (a): optimal discounted cost at time zero, $J_0(s)$, for 500 decisions and zero terminal cost.}

\section{Part (b): infinite-horizon value iteration}
Starting from $J^{(0)}(s)=0$, update
\[
J^{(n+1)}(s)=\min_{a\in\{0,1\}}\left[c(s,a)+\frac{\gamma}{5}
\sum_{k=1}^{5}J^{(n)}\!\left(\min\{q(s,a)+k,200\}\right)\right].
\]
The index $n$ counts iterations, not decision periods. Let $\mathcal TJ$ denote this entire update applied to a candidate vector $J$. Stop when the Bellman residual $\delta(J)=\max_s|(\mathcal TJ)(s)-J(s)|$ is at most $10^{-10}$. This took \textbf{518 updates}, with final residual $9.913\times10^{-11}$.
\shuttlefigure{b_value_iteration_value.pdf}{Part (b): infinite-horizon optimal cost $J^*(s)$ from value iteration. Its near agreement with Part (a) is tested using a discounted-tail bound.}

\clearpage
\section{Part (c): policy iteration and optimal policy}
Start from always waiting. For a stationary policy $\pi$, form its transition matrix $P_\pi$ and cost vector $c_\pi$, where $P_\pi(s,s')=P(s'\mid s,\pi(s))$ and $c_\pi(s)=c(s,\pi(s))$. Its value vector $J_\pi$ satisfies
\[
(I-\gamma P_\pi)J_\pi=c_\pi.
\]
$I$ is the identity matrix. Solve this linear system, then improve each state's action by comparing current cost plus discounted continuation under $J_\pi$. Repeat until no action changes. Policy iteration took \textbf{six evaluation/improvement rounds}, including the final stability check.

The optimal stationary policy is to \textbf{wait at queues 0--12 and dispatch at queues 13--200}. No threshold restriction was imposed on the search.
\shuttlefigure{c_policy_iteration_policy.pdf}{Part (c): optimal stationary policy. Dispatch begins at 13 passengers, before the shuttle is full.}

\textbf{Why dispatch before 15?} The action costs, including optimal future costs, switch order between 12 and 13:
\begin{center}
\begin{tabular}{rrrl}\toprule
Queue & Wait cost & Dispatch cost & Decision\\\midrule
12 & \WaitCostTwelve & \DispatchCostTwelve & Wait\\
13 & \WaitCostThirteen & \DispatchCostThirteen & Dispatch\\\bottomrule
\end{tabular}\end{center}
Waiting for a full shuttle is feasible but more expensive: exact policy evaluation gives cost \FullShuttleCost{} from an empty queue, versus \OptimalEmptyCost{} for the optimal policy. Representative optimal costs are $J^*(15)=676.670168$ and $J^*(200)=3821.586477$.

\subsection*{Edge cases and their consequences}
\textbf{Empty and underfilled shuttles.} The maximum with zero prevents negative queues. Dispatching at 0 costs 100 extra without changing the next-state distribution, so it is dominated. Dispatch at 15 clears the queue; at 16 it leaves one. Also, $J^*(15)-J^*(0)=100$, since optimal dispatch at 15 and waiting at 0 leave identical queues before arrivals.

\textbf{Station capacity.} At 200, waiting leads to 200 with probability one; dispatch leads to 186--190 with probability $1/5$ each. Repeated capped outcomes must be pooled, not treated as distinct next states with lost probability.

\textbf{The final decision.} Waiting is optimal everywhere at $t=499$: dispatch costs 100 but saves at most $2\times15=30$ immediately. With no terminal penalty, there is no later benefit. The saved final-period policy and terminal value were checked against this prediction; the infinite-horizon rule must not be used at every finite-horizon time.

\clearpage
\section{Quantitative validation and execution}\label{sec:validation}
All \textbf{14 saved validation experiments passed}. Agreement alone is insufficient because solvers share a model; the tests also use analytic modifications, independent tree enumeration and independently coded simulation.

\textbf{Optimality and solver agreement.} Value and policy iteration choose identical actions in all 201 states; their maximum value difference is $1.964\times10^{-9}$. Policy iteration's Bellman residual is $9.095\times10^{-13}$. Since discounting makes the update contractive,
\[
\max_s|J(s)-J^*(s)|\leq\frac{\delta(J)}{1-\gamma},
\]
so the residual implies a value-error bound of $1.819\times10^{-11}$, subject to rounding. The smallest action-cost gap is 3.778, far larger than numerical error. Finite-horizon time-zero values differ from infinite-horizon values by at most $4.942\times10^{-9}$, below the omitted-tail bound $470\gamma^{500}/(1-\gamma)=6.838\times10^{-8}$; 470 bounds any stage cost.

\textbf{Modified problems with known answers.} Each prediction below matched all states:
\begin{center}\small
\begin{tabularx}{\linewidth}{@{}p{.37\linewidth}X@{}}\toprule
Modification & Predicted and observed result\\\midrule
Holding cost set to zero & Always wait; optimal cost zero.\\
Zero capacity and no arrivals & Always wait; $J(s)=2s/(1-\gamma)$; zero analytic error.\\
Free dispatch, capacity 200, no arrivals & Dispatch at nonempty queues; cost zero.\\
Discount set to zero & Always wait; $J(s)=2s$.\\\bottomrule
\end{tabularx}\end{center}
An exhaustive scalar decision/arrival tree with seven states and four periods matched backward induction exactly. This checks the recursion without reusing its vectorized transition arrays.

\textbf{Independent simulation.} For each initial queue, simulate 20,000 trajectories of 300 decisions with seed 1668, coding costs/transitions directly rather than using the solver's array builder:
\begin{center}\small
\begin{tabular}{rrrr}\toprule
Initial queue & Simulated mean & Exact policy cost & Standard error\\\midrule
0 & 576.671469 & 576.670168 & 0.219213\\
15 & 676.814669 & 676.670168 & 0.221520\\
200 & 3821.605190 & 3821.586477 & 0.532280\\\bottomrule
\end{tabular}\end{center}
Standard error is trajectory-cost sample SD divided by $\sqrt{20{,}000}$. Every discrepancy is below one SE; the automated tolerance is four SE plus a 0.001951 omitted-tail bound. Simulation checks policy evaluation; Bellman optimality supplies the optimality evidence.

\textbf{Rebuild the report.} Run \texttt{python make\_report\_figures.py}, then compile \texttt{SHUTTLE\_REPORT.tex}. The builder audits saved arrays and reproduces the full-shuttle comparison without rerunning or overwriting the original solvers.

\textbf{With three additional months:} test cost, capacity, arrival and queue-cap sensitivity; compare overflow penalties and terminal costs. Independently solve the MDP by linear programming and examine near-tied actions at higher precision. Finally, test the policy under correlated demand and explicit shuttle availability/travel times, adding those variables to the state where necessary. These distinguish numerical verification from robustness to modeling assumptions.
\end{document}
